\subsection*{Null and Constant Sequences}

% ---------------------------------------------------------
% Definition --- Constant Sequence
% ---------------------------------------------------------

\begin{definitionbox}{Definition (Constant Sequence)}
\begin{definition}[Constant Sequence]
\label{def:constant-sequence}
Let $(x_n)$ be a real sequence. The sequence $(x_n)$ is
\textbf{constant} exactly when there exists $c \in \mathbb{R}$ such that
$x_n = c$ for every $n \in \mathbb{N}$.
\end{definition}
\end{definitionbox}

\begin{remark*}[Standard quantified statement]
\[
\begin{aligned}
&\text{Let } (x_n) \text{ be a real sequence.}\\
&(x_n) \text{ is constant}
\Longleftrightarrow
\exists c \in \mathbb{R}\;
\forall n \in \mathbb{N}\;(x_n = c).
\end{aligned}
\]
\end{remark*}

\begin{remark*}[Definition predicate reading]
\[
\operatorname{ConstantSequence}(x_n,c)
\coloneqq
\forall n \in \mathbb{N}\;(x_n = c).
\]
\end{remark*}

\begin{remark*}[Negated quantified statement]
\[
\begin{aligned}
&\text{Let } (x_n) \text{ be a real sequence.}\\
&(x_n) \text{ is not constant}
\Longleftrightarrow
\forall c \in \mathbb{R}\;
\exists n \in \mathbb{N}\;(x_n \neq c).
\end{aligned}
\]
\end{remark*}

\begin{remark*}[Negation predicate reading]
\[
\neg\exists c \in \mathbb{R}\;
\operatorname{ConstantSequence}(x_n,c).
\]
\end{remark*}

\begin{remark*}[Failure modes]
Every proposed constant value fails at some index.
\end{remark*}

\begin{remark*}[Failure mode decomposition]
\[
\forall c \in \mathbb{R}\;
\underbrace{\exists n \in \mathbb{N}\;(x_n \neq c)}
_{\text{a term differs from the proposed constant value}}.
\]
\end{remark*}

\begin{remark*}[Interpretation]
A constant sequence has no genuine dependence on the index. The indexing is
still part of the object, but every index is assigned the same real value.
\end{remark*}

\begin{dependencies}
\begin{itemize}
  \item \hyperref[def:sequence]{Sequence}
\end{itemize}
\end{dependencies}

% ---------------------------------------------------------
% Definition --- Null Sequence
% ---------------------------------------------------------

\begin{definitionbox}{Definition (Null Sequence)}
\begin{definition}[Null Sequence]
\label{def:null-sequence}
Let $(x_n)$ be a real sequence. The sequence $(x_n)$ is a
\textbf{null sequence} exactly when for every $\varepsilon > 0$ there exists
$N \in \mathbb{N}$ such that $|x_n| < \varepsilon$ for every $n \ge N$.
\end{definition}
\end{definitionbox}

\begin{remark*}[Standard quantified statement]
\[
\begin{aligned}
&\text{Let } (x_n) \text{ be a real sequence.}\\
&(x_n) \text{ is a null sequence}
\Longleftrightarrow
\forall \varepsilon > 0\;
\exists N \in \mathbb{N}\;
\forall n \ge N\;(|x_n| < \varepsilon).
\end{aligned}
\]
\end{remark*}

\begin{remark*}[Definition predicate reading]
\[
\operatorname{NullSequence}(x_n)
\coloneqq
\forall \varepsilon > 0\;
\exists N \in \mathbb{N}\;
\forall n \ge N\;(|x_n| < \varepsilon).
\]
\end{remark*}

\begin{remark*}[Negated quantified statement]
\[
\begin{aligned}
&\text{Let } (x_n) \text{ be a real sequence.}\\
&(x_n) \text{ is not a null sequence}
\Longleftrightarrow
\exists \varepsilon > 0\;
\forall N \in \mathbb{N}\;
\exists n \ge N\;(|x_n| \ge \varepsilon).
\end{aligned}
\]
\end{remark*}

\begin{remark*}[Negation predicate reading]
\[
\neg\operatorname{NullSequence}(x_n).
\]
\end{remark*}

\begin{remark*}[Failure modes]
Some positive tolerance is missed infinitely far out in the sequence.
\end{remark*}

\begin{remark*}[Failure mode decomposition]
\[
\exists \varepsilon > 0\;
\forall N \in \mathbb{N}\;
\underbrace{\exists n \ge N\;(|x_n| \ge \varepsilon)}
_{\text{the tail is not trapped inside the } \varepsilon\text{-band}}.
\]
\end{remark*}

\begin{remark*}[Interpretation]
A null sequence is a sequence whose tail is eventually contained in every
symmetric neighborhood of zero. The first finitely many terms are irrelevant;
only the eventual behavior matters.
\end{remark*}

\begin{dependencies}
\begin{itemize}
  \item \hyperref[def:sequence]{Sequence}
\end{itemize}
\end{dependencies}

% ---------------------------------------------------------
% Theorem --- Constant Sequence Convergence
% ---------------------------------------------------------

\begin{theorembox}{Theorem (Constant Sequence Convergence)}
\begin{theorem}[Constant Sequence Convergence]
\label{thm:constant-sequence-convergence}
Let $(x_n)$ be a real sequence, and let $c \in \mathbb{R}$. If $x_n = c$
for every $n \in \mathbb{N}$, then $(x_n)$ converges to $c$.

\smallskip
\noindent
\hyperref[prf:constant-sequence-convergence]{\textit{Go to proof.}}
\end{theorem}
\end{theorembox}

\begin{remark*}[Standard quantified statement]
\[
\begin{aligned}
&\text{Let } (x_n) \text{ be a real sequence and } c \in \mathbb{R}.\\
&\bigl(\forall n \in \mathbb{N}\;(x_n = c)\bigr)
\Longrightarrow
\forall \varepsilon > 0\;
\exists N \in \mathbb{N}\;
\forall n \ge N\;(|x_n - c| < \varepsilon).
\end{aligned}
\]
\end{remark*}

\begin{remark*}[Predicate reading]
\[
\operatorname{ConstantSequence}(x_n,c)
\Longrightarrow
\operatorname{ConvergentSequence}(x_n,c).
\]
\end{remark*}

\begin{remark*}[Interpretation]
A constant sequence is already equal to its proposed limit at every index.
Every positive tolerance is therefore satisfied by any tail of the sequence.
\end{remark*}

\begin{dependencies}
\begin{itemize}
  \item \hyperref[def:constant-sequence]{Constant Sequence}
  \item \hyperref[def:convergent-sequence]{Convergent Sequence}
\end{itemize}
\end{dependencies}

% ---------------------------------------------------------
% Theorem --- Zero Sequence is Null
% ---------------------------------------------------------

\begin{theorembox}{Theorem (Zero Sequence is Null)}
\begin{theorem}[Zero Sequence is Null]
\label{thm:zero-sequence-is-null}
Let $(x_n)$ be a real sequence. If $x_n = 0$ for every $n \in \mathbb{N}$,
then $(x_n)$ is a null sequence.

\smallskip
\noindent
\hyperref[prf:zero-sequence-is-null]{\textit{Go to proof.}}
\end{theorem}
\end{theorembox}

\begin{remark*}[Standard quantified statement]
\[
\begin{aligned}
&\text{Let } (x_n) \text{ be a real sequence.}\\
&\bigl(\forall n \in \mathbb{N}\;(x_n = 0)\bigr)
\Longrightarrow
\forall \varepsilon > 0\;
\exists N \in \mathbb{N}\;
\forall n \ge N\;(|x_n| < \varepsilon).
\end{aligned}
\]
\end{remark*}

\begin{remark*}[Predicate reading]
\[
\operatorname{ConstantSequence}(x_n,0)
\Longrightarrow
\operatorname{NullSequence}(x_n).
\]
\end{remark*}

\begin{remark*}[Interpretation]
The zero sequence satisfies every null-sequence tolerance immediately because
the distance from each term to zero is exactly zero.
\end{remark*}

\begin{dependencies}
\begin{itemize}
  \item \hyperref[def:null-sequence]{Null Sequence}
  \item \hyperref[def:constant-sequence]{Constant Sequence}
\end{itemize}
\end{dependencies}

% ---------------------------------------------------------
% Theorem --- Constant Null Sequence
% ---------------------------------------------------------

\begin{theorembox}{Theorem (Constant Null Sequence)}
\begin{theorem}[Constant Null Sequence]
\label{thm:constant-null-sequence}
Let $(x_n)$ be a real sequence, and let $c \in \mathbb{R}$. If $x_n = c$
for every $n \in \mathbb{N}$, then $(x_n)$ is a null sequence if and only if
$c = 0$.

\smallskip
\noindent
\hyperref[prf:constant-null-sequence]{\textit{Go to proof.}}
\end{theorem}
\end{theorembox}

\begin{remark*}[Standard quantified statement]
\[
\begin{aligned}
&\text{Let } (x_n) \text{ be a real sequence and } c \in \mathbb{R}.\\
&\bigl(\forall n \in \mathbb{N}\;(x_n = c)\bigr)
\Longrightarrow
\Bigl(
\bigl[
\forall \varepsilon > 0\;
\exists N \in \mathbb{N}\;
\forall n \ge N\;(|x_n| < \varepsilon)
\bigr]
\Longleftrightarrow c = 0
\Bigr).
\end{aligned}
\]
\end{remark*}

\begin{remark*}[Predicate reading]
\[
\operatorname{ConstantSequence}(x_n,c)
\Longrightarrow
\bigl(\operatorname{NullSequence}(x_n) \Longleftrightarrow c = 0\bigr).
\]
\end{remark*}

\begin{remark*}[Interpretation]
A sequence that never changes can approach zero only when its fixed value is
already zero. Conversely, the fixed value zero gives the zero sequence.
\end{remark*}

\begin{dependencies}
\begin{itemize}
  \item \hyperref[def:constant-sequence]{Constant Sequence}
  \item \hyperref[def:null-sequence]{Null Sequence}
\end{itemize}
\end{dependencies}

% ---------------------------------------------------------
% Definition --- Ultimately Constant Sequence
% ---------------------------------------------------------

\begin{definitionbox}{Definition (Ultimately Constant Sequence)}
\begin{definition}[Ultimately Constant Sequence]
\label{def:ultimately-constant-sequence}
Let $(x_n)$ be a real sequence. The sequence $(x_n)$ is
\textbf{ultimately constant} exactly when there exist $N \in \mathbb{N}$ and
$c \in \mathbb{R}$ such that $x_n = c$ for every $n \ge N$.
\end{definition}
\end{definitionbox}

\begin{remark*}[Standard quantified statement]
\[
\begin{aligned}
&\text{Let } (x_n) \text{ be a real sequence.}\\
&(x_n) \text{ is ultimately constant}
\Longleftrightarrow
\exists N \in \mathbb{N}\;
\exists c \in \mathbb{R}\;
\forall n \ge N\;(x_n = c).
\end{aligned}
\]
\end{remark*}

\begin{remark*}[Definition predicate reading]
\[
\operatorname{UltimatelyConstantSequence}(x_n,c)
\coloneqq
\exists N \in \mathbb{N}\;
\forall n \ge N\;(x_n = c).
\]
\end{remark*}

\begin{remark*}[Negated quantified statement]
\[
\begin{aligned}
&\text{Let } (x_n) \text{ be a real sequence.}\\
&(x_n) \text{ is not ultimately constant}
\Longleftrightarrow
\forall N \in \mathbb{N}\;
\forall c \in \mathbb{R}\;
\exists n \ge N\;(x_n \neq c).
\end{aligned}
\]
\end{remark*}

\begin{remark*}[Negation predicate reading]
\[
\neg\exists c \in \mathbb{R}\;
\operatorname{UltimatelyConstantSequence}(x_n,c).
\]
\end{remark*}

\begin{remark*}[Failure modes]
No tail of the sequence becomes permanently equal to a single real value.
\end{remark*}

\begin{remark*}[Failure mode decomposition]
\[
\forall N \in \mathbb{N}\;
\forall c \in \mathbb{R}\;
\underbrace{\exists n \ge N\;(x_n \neq c)}
_{\text{the proposed tail value fails later in the sequence}}.
\]
\end{remark*}

\begin{remark*}[Interpretation]
An ultimately constant sequence may vary on an initial finite segment, but its
tail becomes fixed. The concept isolates eventual behavior from finite
initial behavior.
\end{remark*}

\begin{dependencies}
\begin{itemize}
  \item \hyperref[def:sequence]{Sequence}
\end{itemize}
\end{dependencies}

% ---------------------------------------------------------
% Theorem --- Difference from the Limit is Null
% ---------------------------------------------------------

\begin{theorembox}{Theorem (Difference from the Limit is Null)}
\begin{theorem}[Difference from the Limit is Null]
\label{thm:difference-from-limit-is-null}
Let $(x_n)$ be a real sequence, and let $L \in \mathbb{R}$. The sequence
$(x_n)$ converges to $L$ if and only if the sequence $(x_n - L)$ is a null
sequence.

\smallskip
\noindent
\hyperref[prf:difference-from-limit-is-null]{\textit{Go to proof.}}
\end{theorem}
\end{theorembox}

\begin{remark*}[Standard quantified statement]
\[
\begin{aligned}
&\text{Let } (x_n) \text{ be a real sequence and } L \in \mathbb{R}.\\
&\bigl[
\forall \varepsilon > 0\;
\exists N \in \mathbb{N}\;
\forall n \ge N\;(|x_n - L| < \varepsilon)
\bigr]\\
&\qquad\Longleftrightarrow\qquad
\bigl[
\forall \varepsilon > 0\;
\exists N \in \mathbb{N}\;
\forall n \ge N\;(|x_n - L| < \varepsilon)
\bigr].
\end{aligned}
\]
\end{remark*}

\begin{remark*}[Predicate reading]
\[
\operatorname{ConvergentSequence}(x_n,L)
\Longleftrightarrow
\operatorname{NullSequence}(x_n - L).
\]
\end{remark*}

\begin{remark*}[Interpretation]
Subtracting the proposed limit translates the convergence problem to the
special case of convergence to zero. Null sequences are therefore the local
normal form for sequence convergence.
\end{remark*}

\begin{dependencies}
\begin{itemize}
  \item \hyperref[def:convergent-sequence]{Convergent Sequence}
  \item \hyperref[def:null-sequence]{Null Sequence}
  \item \hyperref[def:pointwise-difference-sequence]{Pointwise Difference of Sequences}
\end{itemize}
\end{dependencies}

% ---------------------------------------------------------
% Theorem --- Ultimately Constant Sequence Convergence
% ---------------------------------------------------------

\begin{theorembox}{Theorem (Ultimately Constant Sequence Convergence)}
\begin{theorem}[Ultimately Constant Sequence Convergence]
\label{thm:ultimately-constant-sequence-convergence}
Let $(x_n)$ be a real sequence. If there exist $N_0 \in \mathbb{N}$ and
$c \in \mathbb{R}$ such that $x_n = c$ for every $n \ge N_0$, then
$(x_n)$ converges to $c$.

\smallskip
\noindent
\hyperref[prf:ultimately-constant-sequence-convergence]{\textit{Go to proof.}}
\end{theorem}
\end{theorembox}

\begin{remark*}[Standard quantified statement]
\[
\begin{aligned}
&\text{Let } (x_n) \text{ be a real sequence.}\\
&\forall c \in \mathbb{R}\;
\Bigl(
\bigl[
\exists N_0 \in \mathbb{N}\;
\forall n \ge N_0\;(x_n = c)
\bigr]\\
&\qquad\Longrightarrow\qquad
\bigl[
\forall \varepsilon > 0\;
\exists N \in \mathbb{N}\;
\forall n \ge N\;(|x_n - c| < \varepsilon)
\bigr]
\Bigr).
\end{aligned}
\]
\end{remark*}

\begin{remark*}[Predicate reading]
\[
\begin{aligned}
\operatorname{UltimatelyConstantSequence}(x_n,c)
\Longrightarrow
\operatorname{ConvergentSequence}(x_n,c).
\end{aligned}
\]
\end{remark*}

\begin{remark*}[Interpretation]
Once the tail is equal to a fixed value, all sufficiently late terms are
already exactly at the proposed limit.
\end{remark*}

\begin{dependencies}
\begin{itemize}
  \item \hyperref[def:ultimately-constant-sequence]{Ultimately Constant Sequence}
  \item \hyperref[def:convergent-sequence]{Convergent Sequence}
\end{itemize}
\end{dependencies}

% ---------------------------------------------------------
% Theorem --- Constant Implies Ultimately Constant
% ---------------------------------------------------------

\begin{theorembox}{Theorem (Constant Implies Ultimately Constant)}
\begin{theorem}[Constant Implies Ultimately Constant]
\label{thm:constant-implies-ultimately-constant}
Every constant real sequence is ultimately constant.

\smallskip
\noindent
\hyperref[prf:constant-implies-ultimately-constant]{\textit{Go to proof.}}
\end{theorem}
\end{theorembox}

\begin{remark*}[Standard quantified statement]
\[
\begin{aligned}
&\text{Let } (x_n) \text{ be a real sequence.}\\
&\bigl[
\exists c \in \mathbb{R}\;
\forall n \in \mathbb{N}\;(x_n = c)
\bigr]\\
&\qquad\Longrightarrow\qquad
\bigl[
\exists N \in \mathbb{N}\;
\exists c \in \mathbb{R}\;
\forall n \ge N\;(x_n = c)
\bigr].
\end{aligned}
\]
\end{remark*}

\begin{remark*}[Predicate reading]
\[
\operatorname{ConstantSequence}(x_n,c)
\Longrightarrow
\operatorname{UltimatelyConstantSequence}(x_n,c).
\]
\end{remark*}

\begin{remark*}[Interpretation]
A constant sequence has already stabilized from the first index, so it is
automatically stabilized on some tail.
\end{remark*}

\begin{dependencies}
\begin{itemize}
  \item \hyperref[def:constant-sequence]{Constant Sequence}
  \item \hyperref[def:ultimately-constant-sequence]{Ultimately Constant Sequence}
\end{itemize}
\end{dependencies}

% ---------------------------------------------------------
% Theorem --- Ultimately Zero Sequence is Null
% ---------------------------------------------------------

\begin{theorembox}{Theorem (Ultimately Zero Sequence is Null)}
\begin{theorem}[Ultimately Zero Sequence is Null]
\label{thm:ultimately-zero-sequence-is-null}
Let $(x_n)$ be a real sequence. If there exists $N_0 \in \mathbb{N}$ such
that $x_n = 0$ for every $n \ge N_0$, then $(x_n)$ is a null sequence.

\smallskip
\noindent
\hyperref[prf:ultimately-zero-sequence-is-null]{\textit{Go to proof.}}
\end{theorem}
\end{theorembox}

\begin{remark*}[Standard quantified statement]
\[
\begin{aligned}
&\text{Let } (x_n) \text{ be a real sequence.}\\
&\bigl[
\exists N_0 \in \mathbb{N}\;
\forall n \ge N_0\;(x_n = 0)
\bigr]\\
&\qquad\Longrightarrow\qquad
\forall \varepsilon > 0\;
\exists N \in \mathbb{N}\;
\forall n \ge N\;(|x_n| < \varepsilon).
\end{aligned}
\]
\end{remark*}

\begin{remark*}[Predicate reading]
\[
\operatorname{UltimatelyConstantSequence}(x_n,0)
\Longrightarrow
\operatorname{NullSequence}(x_n).
\]
\end{remark*}

\begin{remark*}[Interpretation]
If the tail is identically zero, then every sufficiently late term lies inside
every positive tolerance around zero.
\end{remark*}

\begin{dependencies}
\begin{itemize}
  \item \hyperref[def:ultimately-constant-sequence]{Ultimately Constant Sequence}
  \item \hyperref[def:null-sequence]{Null Sequence}
\end{itemize}
\end{dependencies}

% ---------------------------------------------------------
% Theorem --- Ultimately Constant Null Sequence
% ---------------------------------------------------------

\begin{theorembox}{Theorem (Ultimately Constant Null Sequence)}
\begin{theorem}[Ultimately Constant Null Sequence]
\label{thm:ultimately-constant-null-sequence}
Let $(x_n)$ be a real sequence. Suppose there exist $N_0 \in \mathbb{N}$
and $c \in \mathbb{R}$ such that $x_n = c$ for every $n \ge N_0$. Then
$(x_n)$ is a null sequence if and only if $c = 0$.

\smallskip
\noindent
\hyperref[prf:ultimately-constant-null-sequence]{\textit{Go to proof.}}
\end{theorem}
\end{theorembox}

\begin{remark*}[Standard quantified statement]
\[
\begin{aligned}
&\text{Let } (x_n) \text{ be a real sequence.}\\
&\forall c \in \mathbb{R}\;
\Bigl(
\bigl[
\exists N_0 \in \mathbb{N}\;
\forall n \ge N_0\;(x_n = c)
\bigr]\\
&\qquad\Longrightarrow\qquad
\Bigl(
\bigl[
\forall \varepsilon > 0\;
\exists N \in \mathbb{N}\;
\forall n \ge N\;(|x_n| < \varepsilon)
\bigr]
\Longleftrightarrow c = 0
\Bigr)
\Bigr).
\end{aligned}
\]
\end{remark*}

\begin{remark*}[Predicate reading]
\[
\operatorname{UltimatelyConstantSequence}(x_n,c)
\Longrightarrow
\bigl(\operatorname{NullSequence}(x_n) \Longleftrightarrow c = 0\bigr).
\]
\end{remark*}

\begin{remark*}[Interpretation]
Only the eventual value controls whether an ultimately constant sequence is
null. A finite initial segment cannot prevent or create null behavior.
\end{remark*}

\begin{dependencies}
\begin{itemize}
  \item \hyperref[def:ultimately-constant-sequence]{Ultimately Constant Sequence}
  \item \hyperref[def:null-sequence]{Null Sequence}
\end{itemize}
\end{dependencies}

% ---------------------------------------------------------
% Theorem --- Tail Equality Preserves Convergence
% ---------------------------------------------------------

\begin{theorembox}{Theorem (Tail Equality Preserves Convergence)}
\begin{theorem}[Tail Equality Preserves Convergence]
\label{thm:tail-equality-preserves-convergence}
Let $(x_n)$ and $(y_n)$ be real sequences. If there exists
$N_0 \in \mathbb{N}$ such that $x_n = y_n$ for every $n \ge N_0$, then
for every $L \in \mathbb{R}$, $(x_n)$ converges to $L$ if and only if
$(y_n)$ converges to $L$.

\smallskip
\noindent
\hyperref[prf:tail-equality-preserves-convergence]{\textit{Go to proof.}}
\end{theorem}
\end{theorembox}

\begin{remark*}[Standard quantified statement]
\[
\begin{aligned}
&\text{Let } (x_n) \text{ and } (y_n) \text{ be real sequences.}\\
&\bigl[
\exists N_0 \in \mathbb{N}\;
\forall n \ge N_0\;(x_n = y_n)
\bigr]\\
&\qquad\Longrightarrow\qquad
\forall L \in \mathbb{R}\;
\Bigl(
\bigl[
\forall \varepsilon > 0\;
\exists N_x \in \mathbb{N}\;
\forall n \ge N_x\;(|x_n - L| < \varepsilon)
\bigr]\\
&\qquad\qquad\Longleftrightarrow\qquad
\bigl[
\forall \varepsilon > 0\;
\exists N_y \in \mathbb{N}\;
\forall n \ge N_y\;(|y_n - L| < \varepsilon)
\bigr]
\Bigr).
\end{aligned}
\]
\end{remark*}

\begin{remark*}[Predicate reading]
\[
\operatorname{EventuallyEqual}(x_n,y_n)
\Longrightarrow
\forall L \in \mathbb{R}\;
\bigl(
\operatorname{ConvergentSequence}(x_n,L)
\Longleftrightarrow
\operatorname{ConvergentSequence}(y_n,L)
\bigr).
\]
\end{remark*}

\begin{remark*}[Interpretation]
Convergence is a tail property. Changing finitely many initial terms cannot
change whether a sequence converges to a specified real number.
\end{remark*}

\begin{dependencies}
\begin{itemize}
  \item \hyperref[def:convergent-sequence]{Convergent Sequence}
\end{itemize}
\end{dependencies}

% ---------------------------------------------------------
% Definition --- Sequence Bounded Above
% ---------------------------------------------------------

\begin{definitionbox}{Definition (Sequence Bounded Above)}
\begin{definition}[Sequence Bounded Above]
\label{def:sequence-bounded-above}
Let $(x_n)$ be a real sequence. The sequence $(x_n)$ is
\textbf{bounded above} exactly when there exists $M \in \mathbb{R}$ such
that $x_n \le M$ for every $n \in \mathbb{N}$.
\end{definition}
\end{definitionbox}

\begin{remark*}[Standard quantified statement]
\[
\begin{aligned}
&\text{Let } (x_n) \text{ be a real sequence.}\\
&(x_n) \text{ is bounded above}
\Longleftrightarrow
\exists M \in \mathbb{R}\;
\forall n \in \mathbb{N}\;(x_n \le M).
\end{aligned}
\]
\end{remark*}

\begin{remark*}[Definition predicate reading]
\[
\operatorname{BoundedAboveSequence}(x_n)
\coloneqq
\exists M \in \mathbb{R}\;
\forall n \in \mathbb{N}\;(x_n \le M).
\]
\end{remark*}

\begin{remark*}[Negated quantified statement]
\[
\begin{aligned}
&\text{Let } (x_n) \text{ be a real sequence.}\\
&(x_n) \text{ is not bounded above}
\Longleftrightarrow
\forall M \in \mathbb{R}\;
\exists n \in \mathbb{N}\;(M < x_n).
\end{aligned}
\]
\end{remark*}

\begin{remark*}[Negation predicate reading]
\[
\neg\operatorname{BoundedAboveSequence}(x_n).
\]
\end{remark*}

\begin{remark*}[Failure modes]
Every proposed upper bound is exceeded by some term of the sequence.
\end{remark*}

\begin{remark*}[Failure mode decomposition]
\[
\forall M \in \mathbb{R}\;
\underbrace{\exists n \in \mathbb{N}\;(M < x_n)}
_{\text{a term rises above the proposed ceiling}}.
\]
\end{remark*}

\begin{remark*}[Interpretation]
A bounded-above sequence has a finite ceiling. Failure means no real number
serves as a ceiling for all terms.
\end{remark*}

\begin{dependencies}
\begin{itemize}
  \item \hyperref[def:sequence]{Sequence}
\end{itemize}
\end{dependencies}

% ---------------------------------------------------------
% Definition --- Sequence Bounded Below
% ---------------------------------------------------------

\begin{definitionbox}{Definition (Sequence Bounded Below)}
\begin{definition}[Sequence Bounded Below]
\label{def:sequence-bounded-below}
Let $(x_n)$ be a real sequence. The sequence $(x_n)$ is
\textbf{bounded below} exactly when there exists $m \in \mathbb{R}$ such
that $m \le x_n$ for every $n \in \mathbb{N}$.
\end{definition}
\end{definitionbox}

\begin{remark*}[Standard quantified statement]
\[
\begin{aligned}
&\text{Let } (x_n) \text{ be a real sequence.}\\
&(x_n) \text{ is bounded below}
\Longleftrightarrow
\exists m \in \mathbb{R}\;
\forall n \in \mathbb{N}\;(m \le x_n).
\end{aligned}
\]
\end{remark*}

\begin{remark*}[Definition predicate reading]
\[
\operatorname{BoundedBelowSequence}(x_n)
\coloneqq
\exists m \in \mathbb{R}\;
\forall n \in \mathbb{N}\;(m \le x_n).
\]
\end{remark*}

\begin{remark*}[Negated quantified statement]
\[
\begin{aligned}
&\text{Let } (x_n) \text{ be a real sequence.}\\
&(x_n) \text{ is not bounded below}
\Longleftrightarrow
\forall m \in \mathbb{R}\;
\exists n \in \mathbb{N}\;(x_n < m).
\end{aligned}
\]
\end{remark*}

\begin{remark*}[Negation predicate reading]
\[
\neg\operatorname{BoundedBelowSequence}(x_n).
\]
\end{remark*}

\begin{remark*}[Failure modes]
Every proposed lower bound is undercut by some term of the sequence.
\end{remark*}

\begin{remark*}[Failure mode decomposition]
\[
\forall m \in \mathbb{R}\;
\underbrace{\exists n \in \mathbb{N}\;(x_n < m)}
_{\text{a term falls below the proposed floor}}.
\]
\end{remark*}

\begin{remark*}[Interpretation]
A bounded-below sequence has a finite floor. Failure means no real number
serves as a floor for all terms.
\end{remark*}

\begin{dependencies}
\begin{itemize}
  \item \hyperref[def:sequence]{Sequence}
\end{itemize}
\end{dependencies}

% ---------------------------------------------------------
% Definition --- Bounded Sequence
% ---------------------------------------------------------

\begin{definitionbox}{Definition (Bounded Sequence)}
\begin{definition}[Bounded Sequence]
\label{def:bounded-sequence}
Let $(x_n)$ be a real sequence. The sequence $(x_n)$ is
\textbf{bounded} exactly when there exists $M > 0$ such that
$|x_n| \le M$ for every $n \in \mathbb{N}$.
\end{definition}
\end{definitionbox}

\begin{remark*}[Standard quantified statement]
\[
\begin{aligned}
&\text{Let } (x_n) \text{ be a real sequence.}\\
&(x_n) \text{ is bounded}
\Longleftrightarrow
\exists M > 0\;
\forall n \in \mathbb{N}\;(|x_n| \le M).
\end{aligned}
\]
\end{remark*}

\begin{remark*}[Definition predicate reading]
\[
\operatorname{BoundedSequence}(x_n)
\coloneqq
\exists M > 0\;
\forall n \in \mathbb{N}\;(|x_n| \le M).
\]
\end{remark*}

\begin{remark*}[Negated quantified statement]
\[
\begin{aligned}
&\text{Let } (x_n) \text{ be a real sequence.}\\
&(x_n) \text{ is not bounded}
\Longleftrightarrow
\forall M > 0\;
\exists n \in \mathbb{N}\;(|x_n| > M).
\end{aligned}
\]
\end{remark*}

\begin{remark*}[Negation predicate reading]
\[
\neg\operatorname{BoundedSequence}(x_n).
\]
\end{remark*}

\begin{remark*}[Failure modes]
Every symmetric bound around zero is escaped by some term of the sequence.
\end{remark*}

\begin{remark*}[Failure mode decomposition]
\[
\forall M > 0\;
\underbrace{\exists n \in \mathbb{N}\;(|x_n| > M)}
_{\text{a term escapes the symmetric band of radius } M}.
\]
\end{remark*}

\begin{remark*}[Interpretation]
A bounded sequence is trapped in a finite symmetric band around zero. This is
equivalent to having both a finite ceiling and a finite floor.
\end{remark*}

\begin{dependencies}
\begin{itemize}
  \item \hyperref[def:sequence]{Sequence}
\end{itemize}
\end{dependencies}

% ---------------------------------------------------------
% Theorem --- Eventually Bounded Above Tail Formulation
% ---------------------------------------------------------

\begin{theorembox}{Theorem (Eventually Bounded Above Tail Formulation)}
\begin{theorem}[Eventually Bounded Above Tail Formulation]
\label{thm:eventually-bounded-above-tail-formulation}
Let $(x_n)$ be a real sequence. The sequence $(x_n)$ is bounded above if and
only if there exist $N_0 \in \mathbb{N}$ and $M \in \mathbb{R}$ such that
$x_n \le M$ for every $n \ge N_0$.

\smallskip
\noindent
\hyperref[prf:eventually-bounded-above-tail-formulation]{\textit{Go to proof.}}
\end{theorem}
\end{theorembox}

\begin{remark*}[Standard quantified statement]
\[
\begin{aligned}
&\text{Let } (x_n) \text{ be a real sequence.}\\
&\bigl[
\exists M \in \mathbb{R}\;
\forall n \in \mathbb{N}\;(x_n \le M)
\bigr]\\
&\qquad\Longleftrightarrow\qquad
\bigl[
\exists N_0 \in \mathbb{N}\;
\exists M \in \mathbb{R}\;
\forall n \ge N_0\;(x_n \le M)
\bigr].
\end{aligned}
\]
\end{remark*}

\begin{remark*}[Predicate reading]
\[
\operatorname{BoundedAboveSequence}(x_n)
\Longleftrightarrow
\exists M \in \mathbb{R}\;
\operatorname{Eventually}\bigl(x_n,\; x_n \le M\bigr).
\]
\end{remark*}

\begin{remark*}[Interpretation]
One-sided boundedness is unchanged by ignoring finitely many initial terms,
because finitely many exceptional values can be absorbed into a larger upper
bound.
\end{remark*}

\begin{dependencies}
\begin{itemize}
  \item \hyperref[def:sequence-bounded-above]{Sequence Bounded Above}
\end{itemize}
\end{dependencies}

% ---------------------------------------------------------
% Theorem --- Eventually Bounded Below Tail Formulation
% ---------------------------------------------------------

\begin{theorembox}{Theorem (Eventually Bounded Below Tail Formulation)}
\begin{theorem}[Eventually Bounded Below Tail Formulation]
\label{thm:eventually-bounded-below-tail-formulation}
Let $(x_n)$ be a real sequence. The sequence $(x_n)$ is bounded below if and
only if there exist $N_0 \in \mathbb{N}$ and $m \in \mathbb{R}$ such that
$m \le x_n$ for every $n \ge N_0$.

\smallskip
\noindent
\hyperref[prf:eventually-bounded-below-tail-formulation]{\textit{Go to proof.}}
\end{theorem}
\end{theorembox}

\begin{remark*}[Standard quantified statement]
\[
\begin{aligned}
&\text{Let } (x_n) \text{ be a real sequence.}\\
&\bigl[
\exists m \in \mathbb{R}\;
\forall n \in \mathbb{N}\;(m \le x_n)
\bigr]\\
&\qquad\Longleftrightarrow\qquad
\bigl[
\exists N_0 \in \mathbb{N}\;
\exists m \in \mathbb{R}\;
\forall n \ge N_0\;(m \le x_n)
\bigr].
\end{aligned}
\]
\end{remark*}

\begin{remark*}[Predicate reading]
\[
\operatorname{BoundedBelowSequence}(x_n)
\Longleftrightarrow
\exists m \in \mathbb{R}\;
\operatorname{Eventually}\bigl(x_n,\; m \le x_n\bigr).
\]
\end{remark*}

\begin{remark*}[Interpretation]
One-sided lower boundedness is unchanged by ignoring finitely many initial
terms, because finitely many exceptional values can be absorbed into a lower
bound.
\end{remark*}

\begin{dependencies}
\begin{itemize}
  \item \hyperref[def:sequence-bounded-below]{Sequence Bounded Below}
\end{itemize}
\end{dependencies}

% ---------------------------------------------------------
% Theorem --- Eventually Bounded Tail Formulation
% ---------------------------------------------------------

\begin{theorembox}{Theorem (Eventually Bounded Tail Formulation)}
\begin{theorem}[Eventually Bounded Tail Formulation]
\label{thm:eventually-bounded-tail-formulation}
Let $(x_n)$ be a real sequence. The sequence $(x_n)$ is bounded if and only
if there exist $N_0 \in \mathbb{N}$ and $M > 0$ such that $|x_n| \le M$ for
every $n \ge N_0$.

\smallskip
\noindent
\hyperref[prf:eventually-bounded-tail-formulation]{\textit{Go to proof.}}
\end{theorem}
\end{theorembox}

\begin{remark*}[Standard quantified statement]
\[
\begin{aligned}
&\text{Let } (x_n) \text{ be a real sequence.}\\
&\bigl[
\exists M > 0\;
\forall n \in \mathbb{N}\;(|x_n| \le M)
\bigr]\\
&\qquad\Longleftrightarrow\qquad
\bigl[
\exists N_0 \in \mathbb{N}\;
\exists M > 0\;
\forall n \ge N_0\;(|x_n| \le M)
\bigr].
\end{aligned}
\]
\end{remark*}

\begin{remark*}[Predicate reading]
\[
\operatorname{BoundedSequence}(x_n)
\Longleftrightarrow
\exists M > 0\;
\operatorname{Eventually}\bigl(x_n,\; |x_n| \le M\bigr).
\]
\end{remark*}

\begin{remark*}[Interpretation]
Boundedness is a tail property up to finitely many exceptions. A finite
initial segment can be enclosed by increasing the bound.
\end{remark*}

\begin{dependencies}
\begin{itemize}
  \item \hyperref[def:bounded-sequence]{Bounded Sequence}
\end{itemize}
\end{dependencies}

% ---------------------------------------------------------
% Theorem --- Bounded Sequence If and Only If Bounded Above and Below
% ---------------------------------------------------------

\begin{theorembox}{Theorem (Bounded Sequence If and Only If Bounded Above and Below)}
\begin{theorem}[Bounded Sequence If and Only If Bounded Above and Below]
\label{thm:bounded-sequence-iff-bounded-above-and-below}
Let $(x_n)$ be a real sequence. The sequence $(x_n)$ is bounded if and only
if it is both bounded above and bounded below.

\smallskip
\noindent
\hyperref[prf:bounded-sequence-iff-bounded-above-and-below]{\textit{Go to proof.}}
\end{theorem}
\end{theorembox}

\begin{remark*}[Standard quantified statement]
\[
\begin{aligned}
&\text{Let } (x_n) \text{ be a real sequence.}\\
&\bigl[
\exists K > 0\;
\forall n \in \mathbb{N}\;(|x_n| \le K)
\bigr]\\
&\qquad\Longleftrightarrow\qquad
\bigl[
\exists m \in \mathbb{R}\;
\exists M \in \mathbb{R}\;
\forall n \in \mathbb{N}\;(m \le x_n \land x_n \le M)
\bigr].
\end{aligned}
\]
\end{remark*}

\begin{remark*}[Predicate reading]
\[
\operatorname{BoundedSequence}(x_n)
\Longleftrightarrow
\bigl(
\operatorname{BoundedBelowSequence}(x_n)
\land
\operatorname{BoundedAboveSequence}(x_n)
\bigr).
\]
\end{remark*}

\begin{remark*}[Interpretation]
An absolute-value bound gives a symmetric interval around zero. Having both a
floor and a ceiling gives an interval that may not be symmetric, but it can be
enlarged to a symmetric one.
\end{remark*}

\begin{dependencies}
\begin{itemize}
  \item \hyperref[def:bounded-sequence]{Bounded Sequence}
  \item \hyperref[def:sequence-bounded-above]{Sequence Bounded Above}
  \item \hyperref[def:sequence-bounded-below]{Sequence Bounded Below}
\end{itemize}
\end{dependencies}

% ---------------------------------------------------------
% Theorem --- Absolute Bound Gives Upper and Lower Bounds
% ---------------------------------------------------------

\begin{theorembox}{Theorem (Absolute Bound Gives Upper and Lower Bounds)}
\begin{theorem}[Absolute Bound Gives Upper and Lower Bounds]
\label{thm:absolute-bound-gives-upper-and-lower-bounds}
Let $(x_n)$ be a real sequence, and let $K > 0$. If $|x_n| \le K$ for every
$n \in \mathbb{N}$, then $-K \le x_n \le K$ for every
$n \in \mathbb{N}$.

\smallskip
\noindent
\hyperref[prf:absolute-bound-gives-upper-and-lower-bounds]{\textit{Go to proof.}}
\end{theorem}
\end{theorembox}

\begin{remark*}[Standard quantified statement]
\[
\begin{aligned}
&\text{Let } (x_n) \text{ be a real sequence.}\\
&\forall K > 0\;
\bigl[
\forall n \in \mathbb{N}\;(|x_n| \le K)
\bigr]\\
&\qquad\Longrightarrow\qquad
\forall n \in \mathbb{N}\;(-K \le x_n \land x_n \le K).
\end{aligned}
\]
\end{remark*}

\begin{remark*}[Interpretation]
An absolute-value inequality simultaneously controls the positive and negative
directions. The upper barrier is $K$, and the lower barrier is $-K$.
\end{remark*}

\begin{dependencies}
\begin{itemize}
  \item \hyperref[def:bounded-sequence]{Bounded Sequence}
  \item \hyperref[def:sequence-bounded-above]{Sequence Bounded Above}
  \item \hyperref[def:sequence-bounded-below]{Sequence Bounded Below}
\end{itemize}
\end{dependencies}

% ---------------------------------------------------------
% Theorem --- Upper and Lower Bounds Give an Absolute Bound
% ---------------------------------------------------------

\begin{theorembox}{Theorem (Upper and Lower Bounds Give an Absolute Bound)}
\begin{theorem}[Upper and Lower Bounds Give an Absolute Bound]
\label{thm:upper-and-lower-bounds-give-absolute-bound}
Let $(x_n)$ be a real sequence. If there exist $m,M \in \mathbb{R}$ such
that $m \le x_n \le M$ for every $n \in \mathbb{N}$, then there exists
$K > 0$ such that $|x_n| \le K$ for every $n \in \mathbb{N}$.

\smallskip
\noindent
\hyperref[prf:upper-and-lower-bounds-give-absolute-bound]{\textit{Go to proof.}}
\end{theorem}
\end{theorembox}

\begin{remark*}[Standard quantified statement]
\[
\begin{aligned}
&\text{Let } (x_n) \text{ be a real sequence.}\\
&\bigl[
\exists m \in \mathbb{R}\;
\exists M \in \mathbb{R}\;
\forall n \in \mathbb{N}\;(m \le x_n \land x_n \le M)
\bigr]\\
&\qquad\Longrightarrow\qquad
\bigl[
\exists K > 0\;
\forall n \in \mathbb{N}\;(|x_n| \le K)
\bigr].
\end{aligned}
\]
\end{remark*}

\begin{remark*}[Interpretation]
Two one-sided barriers can always be replaced by one symmetric barrier around
zero. Any positive $K$ at least as large as both $|m|$ and $|M|$ works.
\end{remark*}

\begin{dependencies}
\begin{itemize}
  \item \hyperref[def:sequence-bounded-above]{Sequence Bounded Above}
  \item \hyperref[def:sequence-bounded-below]{Sequence Bounded Below}
  \item \hyperref[def:bounded-sequence]{Bounded Sequence}
\end{itemize}
\end{dependencies}
